\section{The Change This Book Already Shipped}
\label{sec:ch16-the-change-this-book-already-shipped}

\index{breaking change!the one this book made|(}%
Chapter~\ref{ch:satisfiability} ended on a request as the only check that
catches what composition misses. It is, and there is a difficulty in it that
chapter had no reason to raise. The requests I can send are the ones I thought
of. The requests that decide whether a change is safe are the ones somebody
else already shipped, in an application I have never seen, against a schema
they read a year ago.

\index{breaking change!chapter 14 made one}%
Chapter~\ref{ch:null-blast-radius} made such a change. Five fields gave up
their exclamation marks so that a stopped service would cost a page rather
than a response, and I still think the argument in that chapter is right. It is
also a breaking change to every client that had read those fields as
guaranteed, and nothing in this book noticed. The composition wrote a config at
exit zero. The router served it. Every request the book prints still answers.

\subsection{Five Fields, and Nothing Said}

\index{supergraph!the schema a client actually reads}%
The thing a client reads is not any of the four documents the services export.
It is the composed schema the router serves, which
chapter~\ref{ch:composition} pulled out of the router execution config as
\code{engineConfig.graphqlSchema}. Compose the graph as
chapter~\ref{ch:cross-seam-paging} left it, compose it again as
chapter~\ref{ch:null-blast-radius} left it, and put the two composed documents
side by side. Both are 247 lines, and five of those lines differ. This is what
\code{diff} prints for them, the line under \code{<} being the schema before
and the line under \code{>} the schema now:

\begin{minted}[fontsize=\footnotesize]{text}
18c18
<   ): [Node]!
---
>   ): [Node]
48c48
<   ratingCount: Int!
---
>   ratingCount: Int
182c182
<   session: Session!
---
>   session: Session
210c210
<   ratingCount: Int!
---
>   ratingCount: Int
212c212
<   feedbackUrl: String!
---
>   feedbackUrl: String
\end{minted}

\index{breaking change!five fields, six lines}%
Five lines in the composed schema, and six across the four documents that
produced it, because two subgraphs declare \code{Query.nodes} and composition
merges the two declarations into one field. Two of the five are root fields, \code{nodes} and
\code{ratingCount}. Two more are the fields Ratings contributes to a
\code{Session} it does not own, \code{ratingCount} again and
\code{feedbackUrl}. The fifth is the entity a Search document hydrates.

\index{ratingCount@\code{ratingCount}!two fields of that name}%
One field that looks like a sixth is not.
\code{SessionSearchDocument.ratingCount} is still \code{Int!}, because the
Search service owns that column and answers it out of its own projection. Two
fields in this supergraph are called \code{ratingCount} and only one of them
moved.

\index{composition!prints nothing about a changed type}%
Neither composition said anything. Exit zero both times, and the whole of what
either printed is the line chapter~\ref{ch:satisfiability} already showed, the
one naming the file it wrote. There was never a mechanism by which it could
have said more: \code{wgc router compose} is handed one set of documents and
has never seen another, so it is not comparing anything. It is not a diff
that came back empty. It is not a diff.

\subsection{Where the List of Breaking Changes Actually Lives}

\index{breaking change!the specification does not define it}%
The definition is not where it should be. The GraphQL specification does not
define a breaking change. The word appears twice in the
type system section, once about custom scalar specification urls and once in
the deprecation prose chapter~\ref{ch:schema-design} quotes, and in neither
place as a category. There is no list.

\index{breaking change!graphql-js defines the list}%
What everything downstream uses is a TypeScript enum in the reference
implementation. \code{graphql-js} ships \code{findBreakingChanges}, and beside
it a \code{BreakingChangeType} with sixteen members~\autocite{graphqljsbreaking}:

\begin{minted}[fontsize=\footnotesize]{text}
TYPE_REMOVED, TYPE_CHANGED_KIND, TYPE_REMOVED_FROM_UNION,
VALUE_REMOVED_FROM_ENUM, REQUIRED_INPUT_FIELD_ADDED,
IMPLEMENTED_INTERFACE_REMOVED, FIELD_REMOVED, FIELD_CHANGED_KIND,
REQUIRED_ARG_ADDED, ARG_REMOVED, ARG_CHANGED_KIND, DIRECTIVE_REMOVED,
DIRECTIVE_ARG_REMOVED, REQUIRED_DIRECTIVE_ARG_ADDED,
DIRECTIVE_REPEATABLE_REMOVED, DIRECTIVE_LOCATION_REMOVED
\end{minted}

\index{breaking change!a dangerous category beside it}%
A second enum of six sits next to it, \code{DangerousChangeType}, for changes
that break a client only if it made an assumption it was not entitled to: a
value added to an enum, a type added to a union, an optional argument added.
The split is worth keeping, because most of what a team argues about lives in
the second list.

\index{FIELD\_CHANGED\_KIND@\code{FIELD\_CHANGED\_KIND}!what chapter 14 did}%
Chapter~\ref{ch:null-blast-radius}'s change is \code{FIELD\_CHANGED\_KIND}.
The function that decides it walks a field's type, and this is the branch that
answers for a field which used to be non-null, with the comment the file
carries on it:

\begin{minted}[fontsize=\footnotesize]{text}
if (isNonNullType(oldType)) {
  // if they're both non-null, make sure the underlying types are compatible
  return (
    isNonNullType(newType) &&
    isChangeSafeForObjectOrInterfaceField(oldType.ofType, newType.ofType)
  );
}
\end{minted}

\index{nullability!the output rule and the input rule are opposite}%
The old type is non-null and the new one is not, so the conjunction fails and
the change is unsafe. A comment further down the same file says the other
direction plainly: moving from nullable to non-null is safe. That is for an
output field. For an argument or an input field the same file carries the
mirror image, that moving from non-null to nullable is the safe direction. A
field that stops promising breaks a caller. An argument that stops demanding
does not. One \code{!} in two positions, and the rule reverses.

\subsection{And the Tests Moved With It}

\index{breaking change!the tests move with the schema}%
Three things failed to notice that change at the top of this section: the
composition, the router, and every request in this book. There is a fourth,
and it is the part I would not have volunteered. A schema has tests, and the
tests assert the schema. When I softened those five fields, the change that
softened them also rewrote every assertion that named them, because a change
that leaves them alone does not pass and therefore does not ship. That is not
a lapse of discipline. It is the only shape such a change can have.

\index{breaking change!why your own tests cannot tell you}%
Which is why a suite of operations kept by the team that owns the schema
cannot tell that team it has broken somebody. The operations move when the
schema moves, in the same change, by the same hand, so nothing in the suite is
older than the edit, and being older than the edit is the whole of what would
make it evidence. Nobody found out here. The next section is what finding out
looks like when somebody does.
\index{breaking change!the one this book made|)}%
